\documentclass[11pt]{article}
\usepackage[utf8]{inputenc}
\usepackage[T1]{fontenc}
\usepackage{amsmath,amssymb}
\usepackage{booktabs}
\usepackage{graphicx}
\usepackage{hyperref}
\usepackage[margin=2.5cm]{geometry}

\title{Formal Languages --- Group Assignment}
\author{Laura Santamaria Espinosa \and Daniela Giraldo Salas}
\date{\today}

\begin{document}
\maketitle

\section{Algorithms and data structures}
\paragraph{Item 1.} FIRST and FOLLOW follow Section~4.4.2 of the Dragon Book.
Iterative fixpoint over productions. Sets stored as Python \texttt{set}.
For every nonterminal $X$, \texttt{FIRST[X]} is initialized to $\emptyset$
(or $\{t\}$ if $X$ is a terminal); a loop applies the three rules of
Algorithm~4.4.2 until no set grows.

\paragraph{Item 2, LL(1).} The predictive table is built with Algorithm~4.31
(Section~4.4.3): for every $A\to\alpha$, add $A\to\alpha$ to $M[A,a]$ for every
$a\in\mathrm{FIRST}(\alpha)$, and also to $M[A,b]$ for every $b\in\mathrm{FOLLOW}(A)$
if $\varepsilon\in\mathrm{FIRST}(\alpha)$.
The LL(1) test is the one at the beginning of Section~4.4.3: no cell of $M$ may
contain more than one production. The parser is the table-driven non-recursive
parser of Algorithm~4.34 (Section~4.4.4), with an explicit stack.

\paragraph{Item 2, SLR(1).} The canonical collection of LR(0) items is built with
\texttt{CLOSURE} and \texttt{GOTO} (Section~4.6.2), after augmenting the grammar
with a new start symbol $S'$. The SLR(1) table follows Section~4.6.4: shifts from
the transitions of the automaton, reduces $A\to\beta$ on $\mathrm{FOLLOW}(A)$, and
\texttt{accept} on $\$$ for $S'\to S$. The SLR(1) test is: no cell of ACTION has
more than one entry. The parser is the LR parsing algorithm of Section~4.6.3.

\paragraph{Item 3.} The SDD is defined following Sections~5.1.1 (synthesized and
inherited attributes), 5.1.2 (annotated parse trees), 5.2.1 (dependency graphs)
and 5.2.2 (evaluation order). The evaluator implements the dependency-graph
method of Section~5.2.2: for each attribute instance we create one vertex, for
each semantic rule we add one edge from every used attribute to the defined
attribute, and then we compute a topological order. If the graph has a cycle,
the SDD is circular and we stop. The SDD itself is stored as \emph{data} (a
dictionary of semantic rules), so the evaluator is independent of the particular
SDD.

\paragraph{Data structures.}
\begin{itemize}
\item Grammar: \texttt{list} of pairs $(A, \mathrm{rhs})$.
\item FIRST/FOLLOW: \texttt{dict} from nonterminal to \texttt{set} of symbols.
\item LL(1) table: \texttt{dict} from nonterminal to \texttt{dict} from terminal to list of productions.
\item LR(0) automaton: \texttt{list} of \texttt{frozenset} of items; transitions in a \texttt{dict}.
\item SLR(1) table: ACTION \texttt{dict} from $(state, symbol)$ to shift/reduce/accept; GOTO \texttt{dict} from $(state, nonterminal)$ to state.
\item Parse tree: nodes with symbol, production, children and an attribute dictionary.
\item Dependency graph: adjacency \texttt{dict} and in-degree \texttt{dict} (Kahn's algorithm).
\end{itemize}

\section{Item 1 and 2: own grammar}

Grammar:
\begin{align*}
S &\to A \\
A &\to BA \mid \varepsilon \\
B &\to 0 \mid 1 \mid 2
\end{align*}

\paragraph{FIRST and FOLLOW} (computed by the program):
\begin{verbatim}
First(A) = {0, 1, 2, e}
Follow(A) = {$}
First(B) = {0, 1, 2}
Follow(B) = {0, 1, 2, $}
First(S) = {0, 1, 2, e}
Follow(S) = {$}
\end{verbatim}

\paragraph{LL(1) table} (computed by the program):
\begin{center}
\begin{tabular}{c|cccc}
 & 0 & 1 & 2 & \$ \\
\hline
S & $S\to A$ & $S\to A$ & $S\to A$ & $S\to A$ \\
A & $A\to BA$ & $A\to BA$ & $A\to BA$ & $A\to\varepsilon$ \\
B & $B\to 0$ & $B\to 1$ & $B\to 2$ & \\
\end{tabular}
\end{center}
The grammar has no LL(1) conflicts; it is also SLR(1).

\paragraph{Step-by-step parse of} \texttt{0121} with the LL(1) parser:
\begin{center}
\begin{tabular}{lll}
Stack & Input & Action \\
\hline
$S\$ $ & $0121\$$ & $S\to A$ \\
$A\$ $ & $0121\$$ & $A\to BA$ \\
$BA\$ $ & $0121\$$ & $B\to 0$ \\
$0A\$ $ & $0121\$$ & match 0 \\
$A\$ $ & $121\$$ & $A\to BA$ \\
$BA\$ $ & $121\$$ & $B\to 1$ \\
$1A\$ $ & $121\$$ & match 1 \\
$A\$ $ & $21\$$ & $A\to BA$ \\
$BA\$ $ & $21\$$ & $B\to 2$ \\
$2A\$ $ & $21\$$ & match 2 \\
$A\$ $ & $1\$$ & $A\to BA$ \\
$BA\$ $ & $1\$$ & $B\to 1$ \\
$1A\$ $ & $1\$$ & match 1 \\
$A\$ $ & $\$$ & $A\to\varepsilon$ \\
$\$$ & $\$$ & accept \\
\end{tabular}
\end{center}

\section{Item 3: SDD}

\paragraph{Grammar.} The grammar above (3 nonterminals, 6 productions, $A$
recursive, LL(1) and SLR(1)).

\paragraph{Attributes.}
\begin{center}
\begin{tabular}{lllll}
\toprule
Symbol & Attribute & Kind & Type & Meaning \\
\midrule
$S,A$ & len & synthesized & int & length of derived substring \\
$A$ & pos & inherited & int & starting position of the substring \\
$B$ & pos & inherited & int & position of the digit \\
$B$ & val & synthesized & int & value of the digit \\
$S,A$ & sum & synthesized & int & weighted sum $\sum\, d_i \cdot (pos_i+1)$ \\
$0,1,2$ & lexval & synthesized & char & digit read by the scanner \\
\bottomrule
\end{tabular}
\end{center}

\paragraph{Semantic rules.}
\begin{center}
\begin{tabular}{ll}
\toprule
Production & Semantic rules \\
\midrule
$S\to A$ & $A.pos=0$;\; $S.len=A.len$;\; $S.sum=A.sum$ \\
$A\to B A_2$ & $B.pos=A.pos$;\; $A_2.pos=A.pos+1$;\; $A.len=1+A_2.len$;\; $A.sum=B.val\cdot(B.pos+1)+A_2.sum$ \\
$A\to\varepsilon$ & $A.len=0$;\; $A.sum=0$ \\
$B\to 0$ & $B.val=0$ \\
$B\to 1$ & $B.val=1$ \\
$B\to 2$ & $B.val=2$ \\
\bottomrule
\end{tabular}
\end{center}
Subscripts are used only to distinguish body symbols from the head symbol.

\paragraph{Step 3: example by hand.} Input \texttt{0121}. Parse tree and
annotated values:
\begin{verbatim}
S#0 [len=4, sum=12]
+- A#1 [len=4, pos=0, sum=12]
   +- B#2 [pos=0, val=0]
   |  +- 0#3 [lexval=0]
   +- A#4 [len=3, pos=1, sum=12]
      +- B#5 [pos=1, val=1]
      |  +- 1#6 [lexval=1]
      +- A#7 [len=2, pos=2, sum=10]
         +- B#8 [pos=2, val=2]
         |  +- 2#9 [lexval=2]
         +- A#10 [len=1, pos=3, sum=4]
            +- B#11 [pos=3, val=1]
            |  +- 1#12 [lexval=1]
            +- A#13 [len=0, pos=4, sum=0]
\end{verbatim}

Computation:
\begin{itemize}
\item $A\#13$: len = 0, sum = 0.
\item $A\#10$: len = 1, sum = $1\cdot(3+1)+0 = 4$.
\item $A\#7$: len = 2, sum = $2\cdot(2+1)+4 = 10$.
\item $A\#4$: len = 3, sum = $1\cdot(1+1)+10 = 12$.
\item $A\#1$: len = 4, sum = $0\cdot(0+1)+12 = 12$.
\item $S\#0$: len = 4, sum = 12. \textbf{Main value:} $S.sum = 12$.
\end{itemize}

\paragraph{Dependency graph and evaluation order.} One vertex per attribute
instance; an edge $(u,v)$ means $v$ uses $u$. The program computes a topological
order with Kahn's algorithm (Section~5.2.2). For \texttt{0121} a valid order is
the post-order from right to left; the program prints it for the input.

\paragraph{Step 4: classification.} The SDD is \textbf{L-attributed but not
S-attributed}.
It is not S-attributed because it uses the inherited attribute \texttt{pos}.
It is L-attributed because every inherited attribute of a symbol in the body
depends only on attributes of the head or of symbols to its left, and every
synthesized attribute depends only on attributes of the body.

\paragraph{Step 5: implementation.} The SDD is stored as a Python
\texttt{dict} keyed by production, whose values are lists of
\texttt{(defined, used, function)} triples. The evaluator walks the parse tree,
instantiates the rules, builds the dependency graph, runs Kahn's algorithm, and
evaluates the attributes in topological order. If the graph has a cycle it
reports \texttt{circular} and stops.

\paragraph{Step 6: tests.}
\begin{itemize}
\item Shortest string: \texttt{e}. $S.sum=0$.
\item Single digit: \texttt{0}. $S.sum=0$.
\item Recursive case: \texttt{0121}. $S.sum=12$.
\item Deep nesting: \texttt{012012}. $S.sum=25$.
\item Two digits: \texttt{12}. $S.sum=5$.
\item Not in language: \texttt{4}. Output: \texttt{no}.
\item Circular detection: temporarily replacing the rule of $S\to A$ by
$S.sum = S.sum + A.sum$ makes the evaluator print \texttt{circular}.
\end{itemize}

\section{Problems and solutions}
\begin{itemize}
\item \emph{Avoiding the \texttt{list[:-0]} pitfall during SLR reduction.} In
Python, \texttt{lst[:-0]} returns $[]$, not \texttt{lst}. We guard the pop with
\texttt{if rhs:}.
\item \emph{Iterating FIRST/FOLLOW to a fixpoint.} A naive one-pass loop gives
wrong results; we repeat until no set changes.
\item \emph{Distinguishing $e$ ($\varepsilon$) from a terminal in the input.} We
treat the alternative \texttt{e} as the empty body.
\item \emph{Detecting LL(1)/SLR(1) conflicts.} Both tables are stored as lists of
actions; a cell with $>1$ action is a conflict.
\item \emph{Circular SDD detection.} Kahn's algorithm leaves vertices with
positive in-degree, which we report as \texttt{circular}.
\end{itemize}

\end{document}